\documentclass[12pt]{article}

\setlength{\topmargin}{-0.5 in}
\setlength{\textheight}{9.5 in}
\setlength{\oddsidemargin}{-0.4 in}
\setlength{\textwidth}{7.0 in}

\usepackage{amsmath,graphicx,amssymb}
\usepackage{tikz}
\newcommand*\circled[1]{\tikz[baseline=(char.base)]{
		\node[shape=circle,draw,inner sep=2pt] (char) {#1};}}

\pagestyle{empty}
\begin{document}
	
\centerline{QC HW 3
	\hspace{0.4 in}V. Kamat
	\hspace{0.4 in}Spring 2026}

\section*{Group I}

\begin{enumerate}
\item For each of the following expressions, determine whether or not they represent valid qubit states.
\begin{enumerate}
    \item $\frac{1}{2}|0\rangle+\frac{1}{2}|1\rangle$
    \item $e^{-i}|1\rangle$
    \item $2|0\rangle-1|1\rangle$
    \item $\sqrt{\frac{4}{7}}|0\rangle+\sqrt{\frac{3}{7}}|1\rangle$
\end{enumerate}
\item Find a suitable value of $x$ for which $|\psi\rangle = \dfrac{1}{3}|0\rangle+x|1\rangle$ is a valid qubit state. 
\item Find all possible values of $x$ for which $|\psi\rangle = \dfrac{1}{\sqrt{2}}|0\rangle+x|1\rangle$ is a valid qubit state.
\\
\\
\textbf{Hint:} There are infinitely many solutions. Consider using Euler's formula to describe them. 
\item Find complex number(s) $c$ for which the matrix given by 
$
A =
\begin{pmatrix}
\frac{1}{\sqrt{2}} & \frac{1}{\sqrt{2}} \\
\frac{1}{\sqrt{2}} & c
\end{pmatrix}
$ is unitary. 
\item For each of the following qubit states, compute the probability of measuring the outcomes $|0\rangle$ and $|1\rangle$ respectively.
\begin{enumerate}
    \item $\frac{1}{\sqrt{2}}\left(|0\rangle-|1\rangle\right)$
    \item $e^i|0\rangle$
    \item $\sqrt{\frac{1}{3}}|0\rangle+e^{-i}\sqrt{\frac{2}{3}}|1\rangle$
    \item $i \sqrt{p}|0\rangle-\sqrt{1-p}|1\rangle$ for any real number $0\leq p\leq 1$.
\end{enumerate}
\end{enumerate}

\section*{Group II}

\begin{enumerate}
\item In the Bloch sphere representation, a pure qubit state is written as
\[
|\psi\rangle
=
\cos\left(\tfrac{\theta}{2}\right)|0\rangle
+
e^{i\phi}\sin\left(\tfrac{\theta}{2}\right)|1\rangle,
\quad
0 \le \theta \le \pi,\ \ 0 \le \phi < 2\pi.
\]
Using the physical indistinguishability of global phase, explain why values of $\theta$ outside the interval $[0,\pi]$ are not needed.

\item 
Suppose Alice and Bob execute the BB84 protocol and publicly identify those positions where they used the same basis.
They then randomly choose $k$ of these positions to compare publicly in order to check for the presence of an eavesdropper.

Assume an eavesdropper Eve measures each qubit in a randomly chosen basis, and forwards the measurement result to Bob.

\begin{enumerate}
\item Compute the probability that Eve is \textbf{not} detected on a fixed chosen bit.
\item Compute the probability that Eve is not detected on any of the $k$ chosen bits.
\end{enumerate}

\item 
Consider measurement in the $X$-basis $\{|+\rangle,|-\rangle\}$, where
\[
|+\rangle = \frac{|0\rangle+|1\rangle}{\sqrt{2}}, 
\qquad 
|-\rangle = \frac{|0\rangle-|1\rangle}{\sqrt{2}}.
\]
\begin{enumerate}
\item Write down the corresponding (projective) measurement operators
\[
M_{+} = |+\rangle\langle +|,
\qquad
M_{-} = |-\rangle\langle -|.
\]
\item Compute $M_{+}$ and $M_{-}$ explicitly as $2\times 2$ matrices.
\item Verify that both $M_{+}$ and $M_{-}$ are Hermitian.
\item Verify that $M_{+} + M_{-} = I$.
\end{enumerate}

\end{enumerate}

\end{document}
